\documentclass[11pt]{article}
\usepackage[margin=1in]{geometry}
\usepackage{amsmath,amsthm,mathtools}
\usepackage{amssymb}
\usepackage{hyperref}
\usepackage{enumitem}
\usepackage{booktabs}

\newtheorem{theorem}{Theorem}[section]
\newtheorem{lemma}[theorem]{Lemma}
\newtheorem{proposition}[theorem]{Proposition}
\newtheorem{corollary}[theorem]{Corollary}
\theoremstyle{definition}
\newtheorem{definition}[theorem]{Definition}
\newtheorem{remark}[theorem]{Remark}
\newtheorem{conjecture}[theorem]{Conjecture}

\newcommand{\R}{\mathbb{R}}
\newcommand{\C}{\mathbb{C}}
\newcommand{\T}{\mathbb{T}}
\newcommand{\Z}{\mathbb{Z}}
\newcommand{\E}{\mathbb{E}}
\newcommand{\Hdot}[1]{\dot{H}^{#1}}
\newcommand{\cS}{\mathcal{S}}
\newcommand{\cC}{\mathcal{C}}
\newcommand{\cK}{\mathcal{K}}
\DeclareMathOperator{\re}{Re}
\DeclareMathOperator{\im}{Im}
\DeclareMathOperator{\corr}{corr}
\DeclareMathOperator{\sgn}{sgn}

\title{Phase-Conjugation Antisymmetry and Angular\\
Spreading in the 3D Navier--Stokes Nonlinearity}

\author{Ramiro Fontes\\
Quijotic Research\\
\texttt{ramirofontes@gmail.com}}

\date{March 2026}

\begin{document}
\maketitle

\begin{abstract}
We study the angular distribution of $\Hdot{1/2}$ energy
in Fourier space for the 3D incompressible Navier--Stokes
equations on $\T^3$.

\smallskip\noindent\textbf{Exact results.}
We prove that the per-mode nonlinear $\Hdot{1/2}$ energy
transfer $T(k;u)$ is (i)~exactly antisymmetric under
Fourier phase conjugation $\hat u(k)\mapsto\overline{\hat u(k)}$,
(ii)~independent of the Leray projector (pressure drops out),
and (iii)~degree-one homogeneous in $u$.
We also prove that for any random divergence-free
field whose distribution is invariant under
Fourier phase conjugation ($\hat u(k)\mapsto
\overline{\hat u(k)}$), both $\E[T(k;u)]$ and
$\E[\Sigma(u)]$ vanish.

\smallskip\noindent\textbf{Taylor--Green computation.}
For the Taylor--Green vortex, we show by explicit
enumeration of all contributing triads and exact rational
arithmetic that the second-generation nonlinear forcing
at every original mode satisfies
$V(k)\cdot\widetilde R(k)=-1/512<0$, giving
$\cos\alpha(k) = -1$ (perfect anti-alignment).  The
computation exploits the purely imaginary Fourier structure
$\hat u(k)=iV(k)$ and proceeds by direct enumeration of
all 8 modes.

\smallskip\noindent\textbf{Numerical findings.}
Pseudo-spectral simulations at resolutions $128^3$--$256^3$
with six different initial conditions show that developed
Taylor--Green turbulence exhibits persistent angular
spreading ($\Sigma < 0$) at all Sobolev levels
$s\in[0,3/2]$, including the enstrophy ($s=1$).
The spreading is viscosity-independent to four significant
figures, dominated by cross-sector transfer
($|\Sigma_{\rm cross}|/|\Sigma_{\rm within}|\approx 830$),
and robust to 50\% phase destruction.
Triad phases lock to $\pm\pi/2$ in the NS solution ---
a direct consequence of the preserved purely imaginary
Fourier structure --- while random-phase fields and
concentrated jets produce uniform phase distributions.
All three random broadband initial conditions exhibit
enstrophy-level spreading ($\Sigma^{(1)}<0$) at every
snapshot (30/30), providing evidence for generality
beyond the TG symmetry class.
Surgically concentrating a developed NS solution
(boosting $\cC^{(1)}$ to $0.50$ while preserving
phase structure) produces monotone relaxation toward
isotropy at every snapshot (44/44), with the
stretching--spreading ratio $\mathcal{R}\in[0.6,2.2]$
and a near-universal constitutive law
$\Sigma^{(1)}\approx -0.47\,(\cC^{(1)}-1/J)$.

\smallskip\noindent\textbf{Heuristic mechanism and
regularity program.}
We propose a mechanism --- the ``$-i$ rotation'' in the
advection operator --- that explains both the exact TG
result and the numerical observations, and formulate a
conjecture at the enstrophy level ($s=1$) with a concrete
connection to the regularity problem: the NS nonlinearity
simultaneously spreads angular enstrophy and limits
vortex stretching through the same phase mechanism,
suggesting a self-regulating structure that may prevent
blow-up.
\end{abstract}

\tableofcontents
\newpage

%======================================================================
\section{Introduction}
%======================================================================

The 3D incompressible Navier--Stokes equations on
$\T^3 = [0,2\pi]^3$,
\begin{equation}\label{eq:NS}
\partial_t u + (u\cdot\nabla)u + \nabla p = \nu\Delta u,
\quad \nabla\cdot u = 0, \quad u(0) = u_0,
\end{equation}
remain the subject of a central open problem: whether
smooth initial data produce globally smooth solutions
\cite{Fefferman,Leray}.  Partial results constrain
potential singularities \cite{ESS,BKM}, and Tao~\cite{Tao}
has shown that blow-up can occur for averaged versions of
the equations, suggesting that any regularity proof must
use specific structural properties of the nonlinearity.

Standard approaches to regularity treat the nonlinear
term $(u\cdot\nabla)u$ as the adversary.  In this paper
we present analytical and numerical evidence for a
different picture: the nonlinear term acts as a
\emph{regulator} of angular concentration in Fourier
space, spreading energy from concentrated to dispersed
angular directions.

The paper is organized into three layers of decreasing
rigor:
\begin{itemize}[leftmargin=2em]
\item \textbf{Sections~\ref{sec:general}--\ref{sec:TG}:}
  exact algebraic identities and a finite computation,
  all fully proved.
\item \textbf{Section~\ref{sec:results}:} numerical
  findings from pseudo-spectral simulations.
\item \textbf{Sections~\ref{sec:mechanism}--\ref{sec:conjecture}:}
  a heuristic mechanism, a conditional sign criterion,
  and a conjecture.
\end{itemize}
We make every effort to keep these layers clearly
separated.

\subsection{Setup and definitions}

We decompose $S^2$ (the sphere of wavevector directions
$\hat k = k/|k|$) into $J$ approximately equal-area
sectors $\{\Omega_j\}_{j=1}^J$ using a Fibonacci spiral
partition (Section~\ref{sec:method}).

\begin{definition}[Sector energies and concentration]
\label{def:sector}
For a smooth divergence-free field $u$ on $\T^3$
with Fourier coefficients $\hat u(k)$, $k\in\Z^3$:
\begin{align}
E_j &= \sum_{\hat k\in\Omega_j}|k|\,|\hat u(k)|^2,
\qquad E = \sum_j E_j, \qquad p_j = E_j/E,\\
\cC(u) &= \sum_j p_j^2 = \|p\|_{\ell^2}^2
\quad\text{(angular concentration)},\\
\cS(u) &= -\sum_j p_j\log p_j
\quad\text{(angular entropy)}.
\end{align}
Here $E_j$ is the $\Hdot{1/2}$ energy in sector
$\Omega_j$, and $p_j$ is the fraction of total
$\Hdot{1/2}$ energy in that sector.  Note $\cC\in[1/J,1]$
with $\cC=1/J$ for isotropic fields.
\end{definition}

\begin{definition}[Nonlinear energy transfer]
\label{def:transfer}
Let $B(u,u) = -P[(u\cdot\nabla)u]$ denote the
Leray-projected advection, where
$P_{jm}(k) = \delta_{jm}-k_jk_m/|k|^2$ is the Leray
projector in Fourier space.  The Fourier representation is
\begin{equation}\label{eq:B-fourier}
\widehat{B(u,u)}_j(k)
= -P_{jm}(k)\sum_{\substack{p+q=k\\p,q\in\Z^3}}
iq_l\,\hat u_l(p)\,\hat u_m(q),
\end{equation}
with repeated indices summed ($j,l,m\in\{1,2,3\}$).
The per-mode $\Hdot{1/2}$ energy transfer is
\begin{equation}\label{eq:Tk}
T(k;u) = 2\,\re\bigl[|k|\,\hat u(k)^*
\cdot\widehat{B(u,u)}(k)\bigr].
\end{equation}
The angular spreading functional is:
\begin{equation}\label{eq:Sigma}
\Sigma(u) = \frac{d\cC}{dt}\bigg|_{\rm NL}
= \frac{2}{E}\sum_j p_j\Bigl(\sum_{\hat k\in\Omega_j}T(k;u)
- p_j\sum_k T(k;u)\Bigr).
\end{equation}
Spreading corresponds to $\Sigma < 0$ (concentration
decreasing); concentration to $\Sigma > 0$.
\end{definition}

\subsection{Summary of results}

\noindent\textbf{Exact algebraic results}
(Section~\ref{sec:general}):
\begin{enumerate}[leftmargin=2em,label=(A\arabic*)]
\item Phase-conjugation antisymmetry:
  $T(k;\bar u) = -T(k;u)$
  (Theorem~\ref{thm:antisym}).
\item Conjugation-symmetric vanishing: $\E[T(k;u)] = 0$
  and $\E[\Sigma(u)]=0$
  (Theorem~\ref{thm:gaussian}).
\item Pressure invariance
  (Proposition~\ref{prop:pressure}).
\item Homogeneity: $\Sigma(\lambda u)=\lambda\Sigma(u)$
  (Proposition~\ref{prop:homog}).
\end{enumerate}

\noindent\textbf{Taylor--Green computation}
(Section~\ref{sec:TG}):
\begin{enumerate}[leftmargin=2em,label=(B\arabic*)]
\item $V(k)\cdot\widetilde R(k)=-1/512 < 0$ at all
  TG modes (Theorem~\ref{thm:TG}).
\end{enumerate}

\noindent\textbf{Numerical findings}
(Section~\ref{sec:results}):
\begin{enumerate}[leftmargin=2em,label=(N\arabic*)]
\item $\Sigma(u) < 0$ throughout TG-generated turbulence.
\item Spreading rate independent of $\nu$.
\item Phase conjugation flips $\Sigma$ exactly.
\item Random div-free fields: 49\% spreading.
\item Spreading survives 50\% phase randomization.
\item Spreading established within one time step.
\item Cross-sector transfer dominates by $\mathbf{830\times}$.
\item High-$p$ sectors have $\cos\alpha < 0$ (depletion).
\item Multiple ICs: TG spreads; random fields $\Sigma\approx 0$;
  concentrated jet concentrates despite mode-level depletion.
\item Triad phases lock to $\pm\pi/2$ (purely imaginary)
  in the NS solution, with a $4\!:\!1$ bias toward $+\pi/2$.
  Random-phase and jet fields show uniform distributions.
\item $\Sigma/(\cC-1/J)$ independent of $\nu$ to four
  significant figures across $\nu\in[0.002,0.02]$.
\item $\Sigma^{(s)}<0$ for all $s\in[0,3/2]$: the nonlinear
  term spreads angular concentration at every Sobolev level,
  including the enstrophy ($s=1$).
\item Vortex stretching weakly correlated with enstrophy
  concentration ($\corr=+0.14$), but full nonlinear transfer
  produces net spreading ($\corr=-0.43$): cross-sector
  redistribution dominates.
\item $\Sigma^{(1)}/(\cC^{(1)}-1/J)$ independent of $\nu$:
  enstrophy-level constitutive relation is purely nonlinear.
\item All three random broadband ICs have $\Sigma^{(1)}<0$ at
  every snapshot (30/30): enstrophy spreading is more universal
  than $\Hdot{1/2}$ spreading.
\item Mode-level depletion $\kappa_{\rm align}$ undergoes a
  sharp sign transition at $\cC^{(1)}\approx 0.15$: the
  spreading mechanism engages only above a concentration
  threshold.
\item Enstrophy spreading survives 50\% phase randomization
  ($\Sigma^{(1)}<0$ at $\geq 10/11$ concentration levels
  for $f\leq 0.5$).
\item Surgically concentrated fields ($\cC^{(1)}$ boosted
  to $0.20$--$0.50$ while preserving NS phase structure)
  exhibit $\Sigma^{(1)}<0$ at all 44 snapshots ($256^3$),
  with monotone relaxation toward isotropy.
\item The stretching--spreading ratio
  $\mathcal{R}=|\mathcal{W}|/(|\Sigma^{(1)}|
  \cdot\|\omega\|^2)\in[0.6,2.2]$ in the surgical runs,
  \emph{decreasing} with concentration.
\item The constitutive coefficient $g^{(1)} =
  -\Sigma^{(1)}/(\cC^{(1)}-1/J) \approx 0.47\pm 0.04$
  across $\cC^{(1)}\in[0.13, 0.50]$, identifying the
  analytical target for the regularity program.
\end{enumerate}

\noindent\textbf{Heuristic mechanism and conjecture}
(Sections~\ref{sec:mechanism}--\ref{sec:conjecture}):
\begin{enumerate}[leftmargin=2em,label=(H\arabic*)]
\item Cross-sector depletion via the $-i$ rotation
  (Section~\ref{sec:mechanism}).
\item A conditional sign criterion linking the forward
  cascade phase structure to angular spreading
  (Section~\ref{sec:conditional}).
\item Conjecture~\ref{conj:enstrophy}: $\Sigma^{(1)}(u)<0$
  for all smooth NS solutions with $\cC^{(1)} > 1/J$
  (enstrophy-level spreading).
\item A regularity program based on the coupled
  stretching--spreading mechanism
  (Remark~\ref{rem:regularity-final}).
\end{enumerate}

%======================================================================
\section{Exact Algebraic Results}
\label{sec:general}
%======================================================================

\subsection{Phase-conjugation antisymmetry}

\begin{definition}[Phase conjugation]
\label{def:conj}
For a divergence-free field $u$ on $\T^3$, define
$\bar u$ by $\hat{\bar u}(k) = \overline{\hat u(k)}$.
For a real field this gives $\bar u(x)=u(-x)$ (spatial
parity inversion).  All $\Hdot{s}$ norms, sector energies,
and probabilities $p_j$ are unchanged.
\end{definition}

\begin{remark}
This operation is Fourier phase conjugation (equivalently,
spatial parity), \emph{not} dynamical time reversal.
For the Euler equations, the standard time-reversal
involves $u(x,t)\mapsto -u(x,-t)$; for Navier--Stokes,
viscosity already breaks time-reversal symmetry at the
PDE level.  The identity below is therefore an algebraic
oddness property of the trilinear transfer functional
under parity, not a statement about physical reversibility.
\end{remark}

\begin{theorem}[Phase-conjugation antisymmetry]
\label{thm:antisym}
For any smooth divergence-free $u$ on $\T^3$ and
$k\in\Z^3\setminus\{0\}$:
\begin{equation}
T(k;\bar u) = -T(k;u).
\end{equation}
Consequently $\Sigma(\bar u) = -\Sigma(u)$.
\end{theorem}

\begin{proof}
From~\eqref{eq:Tk} and~\eqref{eq:B-fourier}:
\begin{equation}
T(k;u) = 2|k|\re\Bigl[\hat u(k)^*\cdot
\Bigl(-P(k)\sum_{p+q=k}iq_l\hat u_l(p)\hat u(q)
\Bigr)\Bigr].
\end{equation}
Since $\re[\bar w\cdot(-iz)] = \im[\bar w\cdot z]$,
we may write $T(k;u) = 2|k|\,\im[Z(k;u)]$ where
\begin{equation}\label{eq:Z}
Z(k;u) = \sum_{j,l,m}P_{jm}(k)\sum_{p+q=k}
q_l\,\hat u_j(k)^*\hat u_l(p)\hat u_m(q).
\end{equation}
Under $\hat u(k)\to\overline{\hat u(k)}$: the conjugation
reverses $\hat u_j(k)^*\to\hat u_j(k)$ and converts
$\hat u_l(p)\to\overline{\hat u_l(p)}$,
$\hat u_m(q)\to\overline{\hat u_m(q)}$.  Since $P_{jm}(k)$
and $q_l$ are real:
\begin{equation}
Z(k;\bar u) = \overline{Z(k;u)},
\end{equation}
whence $\im[Z(k;\bar u)]=-\im[Z(k;u)]$ and
$T(k;\bar u)=-T(k;u)$.

For $\Sigma$: since $p_j(\bar u) = p_j(u)$ and
$E(\bar u) = E(u)$, the formula~\eqref{eq:Sigma}
gives $\Sigma(\bar u) = -\Sigma(u)$.
\end{proof}

\begin{corollary}[Pure oddness]
\label{cor:odd}
The transfer is purely odd under phase conjugation:
$T^+ := (T(k;u)+T(k;\bar u))/2 \equiv 0$.
\end{corollary}

\subsection{Vanishing for conjugation-symmetric fields}

\begin{theorem}[Vanishing under conjugation symmetry]
\label{thm:gaussian}
Let $u$ be a random smooth divergence-free field on $\T^3$
with $\E[|T(k;u)|]<\infty$ and $\E[|\Sigma(u)|]<\infty$,
whose distribution is invariant under phase conjugation:
$\bar u \stackrel{d}{=} u$, where
$\hat{\bar u}(k) = \overline{\hat u(k)}$.

Then $\E[T(k;u)]=0$ for every $k\neq 0$, and
$\E[\Sigma(u)]=0$.
\end{theorem}

\begin{proof}
By Theorem~\ref{thm:antisym}: $T(k;\bar u) = -T(k;u)$
for every realization.  Since $\bar u\stackrel{d}{=}u$:
\begin{equation}
\E[T(k;u)] = \E[T(k;\bar u)] = \E[-T(k;u)]
= -\E[T(k;u)].
\end{equation}
Hence $\E[T(k;u)]=0$.

For the spreading functional: under phase conjugation,
$E(\bar u) = E(u)$, $p_j(\bar u)=p_j(u)$, and
$T(k;\bar u) = -T(k;u)$.  From the
definition~\eqref{eq:Sigma}, $\Sigma$ is linear in the
values $\{T(k)\}$ with coefficients depending only on
$\{p_j\}$ and $E$.  Since all coefficients are unchanged
under conjugation while each $T(k)$ flips sign:
$\Sigma(\bar u) = -\Sigma(u)$.  By the same distributional
argument, $\E[\Sigma]=0$.
\end{proof}

\begin{remark}[Applicable models]
The hypothesis $\bar u\stackrel{d}{=}u$ holds for any
isotropic Gaussian field, and more generally for any
random field whose Fourier amplitudes are deterministic
and whose phases (on a half-lattice) are i.i.d.\ uniform
or invariant under $\phi\mapsto -\phi$.  It fails for
fields with systematic phase bias, such as developed
NS solutions --- which is precisely why those solutions
can have $\Sigma\neq 0$.
\end{remark}

\subsection{Pressure invariance}

\begin{proposition}[Pressure drops out]
\label{prop:pressure}
$T(k;u) = 2|k|\re[\hat u(k)^*\cdot
(-\widehat{(u\cdot\nabla)u}(k))]$.
That is, $T(k)$ is the same whether computed with the
Leray-projected forcing $\hat B(k)$ or with the raw
advection $-\widehat{(u\cdot\nabla)u}(k)$.
\end{proposition}

\begin{proof}
The Leray projection removes the component parallel
to $k$:
$\hat B(k) = -\widehat{(u\cdot\nabla)u}(k)
+ k\,\beta(k)$
for some scalar $\beta(k)$.  Since $k\cdot\hat u(k)=0$
(divergence-free), the correction $k\,\beta$ is
orthogonal to $\hat u(k)^*$ and drops out of the
inner product.
\end{proof}

\subsection{Factored form and homogeneity}

\begin{proposition}[Factored form]
\label{prop:factored}
\begin{equation}\label{eq:factored}
T(k;u) = 2|k|\,\im\Bigl[
\sum_{p+q=k}(q\cdot\hat u(p))(\hat u(k)^*\cdot\hat u(q))
\Bigr].
\end{equation}
\end{proposition}

\begin{proof}
By Proposition~\ref{prop:pressure}, we may use the raw
advection.  Its $j$-th Fourier component is
$\widehat{(u\cdot\nabla)u}_j(k)
= \sum_{p+q=k}\sum_l iq_l\hat u_l(p)\hat u_j(q)$
(no Leray projector, replacing $P_{jm}$ with $\delta_{jm}$).
Computing $T = 2|k|\im[\cdots]$ and contracting:
$\sum_j \hat u_j(k)^*\hat u_j(q) = \hat u(k)^*\cdot\hat u(q)$
and
$\sum_l q_l\hat u_l(p) = q\cdot\hat u(p)$.
\end{proof}

\begin{proposition}[Homogeneity]
\label{prop:homog}
$\Sigma(\lambda u) = \lambda\,\Sigma(u)$ for $\lambda>0$.
\end{proposition}

\begin{proof}
$T(k;\lambda u) = \lambda^3 T(k;u)$ (cubic in $\hat u$).
$E(\lambda u) = \lambda^2 E(u)$.  $p_j(\lambda u)=p_j(u)$.
From~\eqref{eq:Sigma}: $\Sigma$ picks up $\lambda^3$
from $T$ and $\lambda^{-2}$ from $1/E$, giving
net degree $3-2=1$.
\end{proof}

%======================================================================
\section{The Taylor--Green Computation}
\label{sec:TG}
%======================================================================

\subsection{Fourier structure}

The Taylor--Green vortex
$u_0 = (\sin x_1\cos x_2\cos x_3,\;
-\cos x_1\sin x_2\cos x_3,\;0)$
has Fourier support on $\cK = \{(\pm 1,\pm 1,\pm 1)\}$
(8 wavevectors):
\begin{equation}\label{eq:TG-fourier}
\hat u_1(k) = \frac{-ik_1}{8},\qquad
\hat u_2(k) = \frac{ik_2}{8},\qquad
\hat u_3(k) = 0,\qquad k\in\cK.
\end{equation}
All coefficients are purely imaginary:
$\hat u(k) = iV(k)$ where
$V(k) = (-k_1/8,\,k_2/8,\,0) \in \R^3$.

\begin{lemma}[Properties of $V$]
\label{lem:TG-imag}
For all $k\in\cK$: (i)~$k\cdot V(k) = 0$
(divergence-free, since $k_1^2=k_2^2=1$);
(ii)~$V(-k)=-V(k)$ (reality);
(iii)~$|V(k)|^2 = 1/32$ (uniform amplitude).
\end{lemma}

\subsection{Phase preservation by the nonlinearity}

\begin{proposition}
\label{prop:phase-preserve}
If $\hat u(k)=iV(k)$ with $V(k)\in\R^3$ for all $k$,
then $\hat B(k)=iR(k)$ with $R(k)\in\R^3$.
\end{proposition}

\begin{proof}
The raw advection is
$\widehat{(u\cdot\nabla)u}_j(k')
= \sum_{p+q=k'}\sum_l iq_l\cdot iV_l(p)\cdot iV_j(q)
= i^3\sum_{p+q=k'}\sum_l q_lV_l(p)V_j(q)$.
Since $i^3=-i$, this equals $-iS_j(k')$ with
$S_j(k') = \sum_{p+q=k'}(q\cdot V(p))V_j(q)\in\R$.
The Leray projector $P(k')$ has real entries, and
$B = -P(\text{adv})$, giving
$\hat B_j(k') = -P_{jm}(-iS_m) = iP_{jm}S_m = iR_j$.
\end{proof}

\begin{corollary}\label{cor:persistence}
The purely imaginary structure is preserved by the full NS
evolution: nonlinearity preserves it
(Proposition~\ref{prop:phase-preserve}), and viscosity
$-\nu|k|^2\hat u = -\nu|k|^2\cdot iV$ preserves it by
linearity.
\end{corollary}

\subsection{First generation: no self-transfer}

\begin{proposition}\label{prop:gen1}
$B(u_0,u_0)$ has Fourier support on
$\cK'=\{k\in\Z^3\setminus\{0\}: k_i\in\{0,\pm 2\}\}$,
which is disjoint from $\cK$.  Therefore $T(k;u_0)=0$
for all $k\in\cK$.
\end{proposition}

\begin{proof}
For $p,q\in\cK$: each component of $p+q$ lies in
$\{-2,0,2\}$, never in $\{\pm 1\}$.  So the advection
creates no modes at the original wavevectors.
Since there is no forcing at $k\in\cK$,
$\hat B(k) = 0$ and $T(k) = 0$.
\end{proof}

\subsection{First-generation modes}

Of the 26 nonzero wavevectors in $\cK'$, exactly 8
receive nonzero forcing (the others vanish by
cancellation).  By explicit computation:

\begin{proposition}\label{prop:gen1-modes}
The nonzero first-generation modes are
$\hat B(k') = iR(k')$ with:
\begin{center}
\begin{tabular}{@{}ll@{}}
\toprule
$k'$ & $R(k')$ \\
\midrule
$(2,0,2)$   & $(1/32,\,0,\,-1/32)$ \\
$(2,0,-2)$  & $(1/32,\,0,\,1/32)$ \\
$(-2,0,2)$  & $(-1/32,\,0,\,-1/32)$ \\
$(-2,0,-2)$ & $(-1/32,\,0,\,1/32)$ \\
$(0,2,2)$   & $(0,\,1/32,\,-1/32)$ \\
$(0,2,-2)$  & $(0,\,1/32,\,1/32)$ \\
$(0,-2,2)$  & $(0,\,-1/32,\,-1/32)$ \\
$(0,-2,-2)$ & $(0,\,-1/32,\,1/32)$ \\
\bottomrule
\end{tabular}
\end{center}
\end{proposition}

\begin{proof}
For each $k'\in\cK'$, compute
$S_j(k')=\sum_{p+q=k'}(q\cdot V(p))V_j(q)$ by
enumerating all pairs $p,q\in\cK$ with $p+q=k'$,
then apply $R = P(k')S$.

We show the computation for $k'=(2,0,2)$.  The
constraint $p+q=(2,0,2)$ with $p,q\in\{\pm 1\}^3$
forces $p_1=q_1=1$, $p_3=q_3=1$, and $p_2=-q_2$.
The two contributing pairs are
$(p,q)\in\{((1,s,1),(1,-s,1)):s=\pm 1\}$.

For $s=1$: $q\cdot V(p) = (1,-1,1)\cdot(-1/8,1/8,0)
= -1/4$.  $V(q) = V(1,-1,1) = (-1/8,-1/8,0)$.
Contribution: $(-1/4)(-1/8,-1/8,0) = (1/32,1/32,0)$.

For $s=-1$: $q\cdot V(p) = (1,1,1)\cdot(-1/8,-1/8,0)
= -1/4$.  $V(q) = V(1,1,1) = (-1/8,1/8,0)$.
Contribution: $(-1/4)(-1/8,1/8,0) = (1/32,-1/32,0)$.

Sum: $S(2,0,2) = (1/16, 0, 0)$.
Leray projection: $|k'|^2=8$, $k'\cdot S = 2/16 = 1/8$.
$R = S - k'(k'\cdot S)/|k'|^2
= (1/16,0,0) - (2,0,2)(1/8)/8
= (1/16-1/32,\; 0,\; -1/32)
= (1/32, 0, -1/32)$.

The remaining 7 modes are computed by the same procedure
and verified by the companion script
\texttt{tg\_exact\_symbolic.py} using exact rational
arithmetic.  All 8 values appear in the table in
Proposition~\ref{prop:gen1-modes}.

A companion script (\texttt{tg\_exact\_symbolic.py})
verifies all 8 modes using Python's \texttt{Fraction}
arithmetic with zero floating-point error.
\end{proof}

\subsection{Second generation: exact anti-alignment}

After one time step $\Delta t$, the field has modes
at both $\cK$ (amplitude $O(1)$) and $\cK'$ (amplitude
$O(\Delta t)$).  Second-generation forcing at $k\in\cK$
comes from triads with one leg in $\cK'$ and one in $\cK$.

\begin{definition}[Second-generation forcing]
\label{def:gen2}
For $k\in\cK$, define:
\begin{equation}\label{eq:Rtilde}
\widetilde R_j(k) = \sum_mP_{jm}(k)\sum_{\substack{
q\in\cK\\k-q\in\cK'}}
\bigl[(q\cdot R(k\!-\!q))V_m(q)
+((k\!-\!q)\cdot V(q))R_m(k\!-\!q)\bigr],
\end{equation}
where the first term is the Type-A contribution
(first-generation mode advects, original is advected)
and the second is Type-B (original advects,
first-generation mode is advected).
The second-generation forcing is
$\hat B^{(2)}(k) = i\Delta t\,\widetilde R(k)$.
\end{definition}

\begin{theorem}[Taylor--Green anti-alignment]
\label{thm:TG}
For all $k\in\cK$:
\begin{equation}
\widetilde R(k) = -\frac{1}{16}V(k),
\end{equation}
and therefore:
\begin{equation}
V(k)\cdot\widetilde R(k) = -\frac{|V(k)|^2}{16}
= -\frac{1}{512} < 0.
\end{equation}
The forcing is exactly anti-parallel to the velocity:
$\cos\alpha(k) = -1$ at every TG mode.
\end{theorem}

\begin{proof}
The computation is finite: for each $k\in\cK$, we
enumerate all $q\in\cK$ with $k-q$ among the 8 nonzero
first-generation modes (Proposition~\ref{prop:gen1-modes}),
evaluate the Type-A and Type-B contributions
in~\eqref{eq:Rtilde}, apply the Leray projector, and sum.
We give the full computation for $k_0=(1,1,1)$ and
tabulate all 8 results.

\smallskip\noindent\textbf{Computation at $k=(1,1,1)$.}
The constraint $q\in\cK$, $k-q\in\cK'$ with $k-q$
among the 8 nonzero first-generation modes yields
exactly two contributing triads:
$q=(-1,1,-1)$ with $p'=(2,0,2)$, and
$q=(1,-1,-1)$ with $p'=(0,2,2)$.
(All other $q\in\cK$ give $k-q$ outside the nonzero
first-generation set.)

For $q=(-1,1,-1)$, $p'=(2,0,2)$:
Type~A: $(q\cdot R(p'))V(q)
= ((-1,1,-1)\cdot(1/32,0,-1/32))\,V(-1,1,-1)
= (-1/32+1/32)\,V(-1,1,-1) = 0$.
Type~B: $(p'\cdot V(q))\,R(p')
= ((2,0,2)\cdot(1/8,1/8,0))\,(1/32,0,-1/32)
= (1/4)(1/32,0,-1/32)
= (1/128,0,-1/128)$.

For $q=(1,-1,-1)$, $p'=(0,2,2)$:
Type~A: $(q\cdot R(p'))V(q)
= ((1,-1,-1)\cdot(0,1/32,-1/32))\,V(1,-1,-1)
= (-1/32+1/32)\,V(1,-1,-1) = 0$.
Type~B: $(p'\cdot V(q))\,R(p')
= ((0,2,2)\cdot(-1/8,-1/8,0))\,(0,1/32,-1/32)
= (-1/4)(0,1/32,-1/32)
= (0,-1/128,1/128)$.

Sum before Leray projection:
$(1/128,-1/128,-1/128+1/128)=(1/128,-1/128,0)$.
Leray projection at $k=(1,1,1)$:
$|k|^2=3$, $k\cdot(1/128,-1/128,0)=0$, so
$\widetilde R(1,1,1)=(1/128,-1/128,0)$.

Check: $-V(1,1,1)/16 = -(-1/8,1/8,0)/16
= (1/128,-1/128,0)$. \checkmark

Inner product: $V\cdot\widetilde R
= (-1/8)(1/128)+(1/8)(-1/128) = -1/512$. \checkmark

Cosine: $(V\cdot\widetilde R)^2 = 1/262144
= |V|^2|\widetilde R|^2$, confirming
$\cos^2\alpha=1$, hence $\cos\alpha=-1$.

\smallskip\noindent\textbf{All 8 modes.}
The same enumeration procedure applied to each
$k\in\cK$ gives:
\begin{center}
\begin{tabular}{@{}lll@{}}
\toprule
$k$ & $\widetilde R(k)$ & $-V(k)/16$ \\
\midrule
$(1,1,1)$    & $(1/128,\,-1/128,\,0)$ & $(1/128,\,-1/128,\,0)$ \\
$(1,1,-1)$   & $(1/128,\,-1/128,\,0)$ & $(1/128,\,-1/128,\,0)$ \\
$(1,-1,1)$   & $(1/128,\,1/128,\,0)$  & $(1/128,\,1/128,\,0)$  \\
$(1,-1,-1)$  & $(1/128,\,1/128,\,0)$  & $(1/128,\,1/128,\,0)$  \\
$(-1,1,1)$   & $(-1/128,\,-1/128,\,0)$ & $(-1/128,\,-1/128,\,0)$ \\
$(-1,1,-1)$  & $(-1/128,\,-1/128,\,0)$ & $(-1/128,\,-1/128,\,0)$ \\
$(-1,-1,1)$  & $(-1/128,\,1/128,\,0)$ & $(-1/128,\,1/128,\,0)$ \\
$(-1,-1,-1)$ & $(-1/128,\,1/128,\,0)$ & $(-1/128,\,1/128,\,0)$ \\
\bottomrule
\end{tabular}
\end{center}
In every case $\widetilde R(k)=-V(k)/16$, giving
$V\cdot\widetilde R = -1/512$ and $\cos\alpha=-1$.

The companion script \texttt{tg\_exact\_symbolic.py}
reproduces the full 8-mode enumeration using Python's
\texttt{Fraction} arithmetic with zero floating-point
error, independently verifying every entry in both
tables (Proposition~\ref{prop:gen1-modes} and above).
\end{proof}

\begin{remark}[The $i^2=-1$ mechanism]
\label{rem:i-squared}
The anti-alignment has a transparent algebraic origin.
The TG coefficients carry one factor of $i$
($\hat u = iV$).  The first-generation forcing picks up
$i^3 = -i$ from two coefficients and the advection
operator's $iq_l$, then the sign in $B=-P(\text{adv})$
restores $+i$ (Proposition~\ref{prop:phase-preserve}).
When first-generation modes feed back through a second
advection, another factor of $i$ accumulates: the total
phase for two-step feedback is $i\cdot i = -1$ relative
to the original velocity.
\end{remark}

\begin{remark}[Persistence]
By Corollary~\ref{cor:persistence}, $\hat u(k,t) = iV(k,t)$
for all $t\geq 0$.  This gives $\cos^2\alpha = 1$ at
every mode (the inner product $V\cdot\widetilde R$ is
always real).  The sign $\cos\alpha = -1$ at the original
TG modes is proved at leading order in $\Delta t$ by
Theorem~\ref{thm:TG}; numerical simulations confirm it
persists throughout the evolution.
\end{remark}

%======================================================================
\section{Numerical Methods}
\label{sec:method}
%======================================================================

\subsection{Solver}

We solve~\eqref{eq:NS} in Fourier space on the
$N^3$ grid ($N=64,128$) using an integrating-factor
Runge--Kutta scheme (IF-RK4).  The linear viscous
operator $-\nu|k|^2$ is integrated exactly via the
factor $e^{-\nu|k|^2\Delta t}$, and the nonlinear term
$\hat B(u,u)$ is advanced by classical RK4 in the
integrating-factor frame.  Aliasing is removed by the
$2/3$ rule (zeroing modes with $|k|>N/3$).  The time
step is $\Delta t = 10^{-3}$ for all runs reported here;
stability was verified by halving $\Delta t$ and checking
that all diagnostics change by less than $10^{-6}$.

\subsection{Initial conditions}

We test six initial conditions spanning different symmetry
classes and concentration levels:
\begin{enumerate}[leftmargin=2em]
\item \textbf{Taylor--Green} (TG):
  $u_0 = (\sin x_1\cos x_2\cos x_3,\,
  -\cos x_1\sin x_2\cos x_3,\,0)$.
  Purely imaginary Fourier structure; 8 active modes.
\item \textbf{ABC flow}:
  $u_0 = (A\sin x_3+C\cos x_2,\,B\sin x_1+A\cos x_3,\,
  C\sin x_2+B\cos x_1)$ with $(A,B,C)=(1,0.8,0.6)$.
  A Beltrami field ($\nabla\times u = u$); 12 active modes.
\item \textbf{Random broadband} (3 seeds):
  isotropic energy spectrum $\propto k^4e^{-k^2/32}$
  peaked at $|k|\approx 4$, independent random phases,
  Leray-projected.  Concentration $\cC\approx 0.07$
  (nearly isotropic).
\item \textbf{Concentrated jet}:
  energy in a narrow cone (half-angle $0.3$\,rad) around
  the $x_3$-axis at $|k|\approx 3$, random phases,
  Leray-projected.  Concentration $\cC\approx 0.48$
  (strongly anisotropic).
\end{enumerate}
Each is evolved to $t=5$ ($5{,}000$ time steps) with
diagnostics recorded every $\Delta t=0.5$.

\subsection{Angular sectors}

We use a Fibonacci spiral partition of $S^2$ into $J$
equal-area sectors with centers
$\hat n_j = (\sin\vartheta_j\cos\psi_j,
\sin\vartheta_j\sin\psi_j,\cos\vartheta_j)$,
$\vartheta_j = \arccos(1-2(j+1/2)/J)$,
$\psi_j = 2\pi j/\varphi$, where
$\varphi = (1+\sqrt{5})/2$ is the golden ratio.
Each wavevector $k\neq 0$ is assigned to its nearest
sector center.  Results are reported for
$J\in\{16,32,64,128\}$; the spreading sign $\Sigma<0$
is independent of $J$.

\subsection{Diagnostics}

\begin{definition}[Phase alignment]
\label{def:phase-align}
For a mode $k$ with $\hat u(k)\neq 0$ and
$\hat B(k)\neq 0$:
\begin{equation}
\cos\alpha(k) = \frac{\re[\hat u(k)^*\cdot\hat B(k)]}
{|\hat u(k)|\,|\hat B(k)|}.
\end{equation}
Since $T(k) = 2|k||\hat u||\hat B|\cos\alpha$:
$\cos\alpha>0$ means the forcing amplifies mode $k$,
$\cos\alpha<0$ means it depletes.
\end{definition}

At each diagnostic snapshot we compute: sector energies
$E_j$ and probabilities $p_j$; nonlinear forcing
$\hat B(k)$; per-mode transfer $T(k)$ via~\eqref{eq:Tk};
per-sector transfer $dE_j/dt|_{\rm NL}
=\sum_{\hat k\in\Omega_j}T(k)$;
spreading functional $\Sigma$ via~\eqref{eq:Sigma};
Pearson correlation $\corr(p_j, dE_j|_{\rm NL})$;
per-sector mean phase alignment
$\langle\cos\alpha\rangle_j$; and the within-sector /
cross-sector transfer decomposition (described below).

\subsection{Cross-sector decomposition}

The nonlinear forcing at mode $k\in\Omega_j$ is
decomposed as $\hat B(k) = \hat B^{\rm within}(k)
+ \hat B^{\rm cross}(k)$, where
$\hat B^{\rm within}(k)$ includes only triads $p+q=k$
with the advecting velocity $\hat u(p)$ in the same
sector as $k$ ($\hat p\in\Omega_j$), and
$\hat B^{\rm cross}(k)$ includes the rest.  In
practice, we compute $\hat B^{\rm within}$ by masking
the advecting field to sector $j$ before evaluating
the pseudospectral convolution, then subtracting from
the full $\hat B$.  The induced spreading decomposition
is $\Sigma = \Sigma_{\rm within}+\Sigma_{\rm cross}$,
where $\Sigma_{\rm within}$ (resp.\ $\Sigma_{\rm cross}$)
is computed from~\eqref{eq:Sigma} using only the
within-sector (resp.\ cross-sector) transfers
$T^{\rm within}(k) = 2|k|\re[\hat u^*\cdot
\hat B^{\rm within}]$ (resp.\ $T^{\rm cross}$).

\subsection{Phase surgery}

We test the role of Fourier phase structure by modifying
a developed NS solution:
full phase randomization (random i.i.d.\ uniform
phases, Leray re-projected);
partial randomization (fraction $f$ of modes
randomized);
phase conjugation ($\hat u\mapsto\overline{\hat u}$).
Each modified field is Leray-projected before computing
$\Sigma$.

\subsection{Validation}

The solver conserves $L^2$ energy to machine precision
in the inviscid limit ($\nu=0$, short time).  The TG
solution at $t=0$ reproduces the exact Fourier coefficients
from~\eqref{eq:TG-fourier}, and the first-generation
forcing matches Proposition~\ref{prop:gen1-modes} to 15
digits.  The antisymmetry $T(k;\bar u)=-T(k;u)$ holds to
$10^{-16}$ relative error for all developed solutions.

%======================================================================
\section{Numerical Results}
\label{sec:results}
%======================================================================

All results below use $N=128$, $\nu=0.01$, $J=32$
unless otherwise stated, with the TG vortex evolved to
$t=5$.

\subsection{Persistent spreading along TG orbit (N1)}

Along the TG trajectory at $N=128$, $\nu=0.01$:
at all 1{,}001 diagnostic snapshots from $t=0.01$ to
$t=20$, $\Sigma < 0$ with
$\corr(p_j,dE_j|_{\rm NL}) = -0.56\pm 0.03$.

\subsection{Viscosity independence (N2)}

\begin{table}[h]
\centering
\caption{Spreading diagnostics at $t=5$, $N=128$, $J=64$,
TG initial data.  The nonlinear correlation is independent
of $\nu$ to three significant figures.}
\label{tab:nu}
\begin{tabular}{@{}ccr@{}}
\toprule
$\nu$ & $\corr(p_j,dE_j|_{\rm NL})$
  & $|{\rm NL}/{\rm viscous}|$ \\
\midrule
$10^{-2}$          & $-0.562$ & $610\times$ \\
$5\times 10^{-3}$  & $-0.562$ & $2{,}640\times$ \\
$2\times 10^{-3}$  & $-0.563$ & $13{,}200\times$ \\
\bottomrule
\end{tabular}
\end{table}

\subsection{Phase-conjugation antisymmetry (N3)}

\begin{center}
\begin{tabular}{@{}lcc@{}}
\toprule
& $\corr$ & $\Sigma$ \\
\midrule
NS solution     & $-0.6022$ & $-1.199\times 10^{-3}$ \\
Conjugated      & $+0.6022$ & $+1.199\times 10^{-3}$ \\
\bottomrule
\end{tabular}
\end{center}
Confirmation of Theorem~\ref{thm:antisym} to machine
precision.

\subsection{Non-universality for random fields (N4)}

Random div-free fields with the same energy spectrum:
22/45 spread (49\%), $\corr\in[-0.7,+0.9]$,
mean $+0.07$.  Consistent with
Theorem~\ref{thm:gaussian}: conjugation-symmetric fields
have $\E[\Sigma]=0$.

\subsection{Robustness (N5)}

Starting from a developed NS solution and randomizing
a fraction $f$ of Fourier phases:
\begin{center}
\begin{tabular}{@{}ccc@{}}
\toprule
$f$ & mean $\corr$ & spread \\
\midrule
$0\%$ & $-0.60$ & 5/5 \\
$30\%$ & $-0.29$ & 5/5 \\
$50\%$ & $-0.34$ & 5/5 \\
$70\%$ & $-0.04$ & 3/5 \\
$100\%$ & $+0.16$ & 1/5 \\
\bottomrule
\end{tabular}
\end{center}

\subsection{Onset (N6)}

$\corr = -0.61$ at $t = 0.01$ (one time step from TG),
stable to $t = 20$.

\subsection{Cross-sector dominance (N7)}

$\Sigma_{\rm within} = -1.4\times 10^{-6}$,
$\Sigma_{\rm cross} = -1.2\times 10^{-3}$
(ratio $\mathbf{830\times}$).  Within-sector self-interaction
\emph{concentrates} ($\corr(p_j, dE^{\rm within}_j)=+0.74$);
cross-sector interaction \emph{spreads}
($\corr(p_j, dE^{\rm cross}_j) = -0.69$).

\subsection{Phase alignment (N8)}

$\corr(p_j,\langle\cos\alpha\rangle_j) = -1.000$.
For this TG-based solution at $t=5$, the four high-$p$
sectors have $\cos\alpha = -1.000$ and low-$p$ sectors
have $\cos\alpha = +1.000$ (to numerical precision).
Theorem~\ref{thm:TG} proves exact anti-alignment at
leading order in $\Delta t$; the persistence to $t=5$
is a numerical observation, consistent with the preserved
imaginary structure (Corollary~\ref{cor:persistence})
but not directly proved by Theorem~\ref{thm:TG}.
For non-TG initial data, partial anti-alignment
($\cos\alpha<0$ on average in high-$p$ sectors)
is expected.

\subsection{Sector count independence}

\begin{center}
\begin{tabular}{@{}ccc@{}}
\toprule
$J$ & $\corr$ & $\Sigma<0$? \\
\midrule
$16$  & $-0.93$ & Yes \\
$32$  & $-0.61$ & Yes \\
$64$  & $-0.56$ & Yes \\
$128$ & $-0.54$ & Yes \\
\bottomrule
\end{tabular}
\end{center}

\subsection{Multiple initial conditions (N9)}
\label{sec:multi-ic}

To test whether spreading is specific to TG symmetry,
we repeat the spreading diagnostic for all six initial
conditions (Section~\ref{sec:method}), each evolved to
$t=5$ at $N=128$, $\nu=0.01$, $J=32$.

\begin{table}[h]
\centering
\caption{Spreading across initial conditions.
$n_{\rm spread}$: snapshots with $\Sigma<0$ out of 11.}
\label{tab:multi-ic}
\begin{tabular}{@{}lcccl@{}}
\toprule
IC & $\Sigma(t\!=\!5)$ & $\corr$ & $n_{\rm spread}$
  & status \\
\midrule
Taylor--Green    & $-1.20\times 10^{-3}$ & $-0.69$ & 11/11
  & spreads \\
ABC flow         & $\approx 0$     & $+1.00$ & 2/11
  & degenerate \\
Random (seed 42) & $+9.1\times 10^{-12}$ & $+0.13$ & 0/11
  & concentrates \\
Random (seed 99) & $-1.3\times 10^{-12}$ & $-0.18$ & 11/11
  & spreads \\
Random (seed 7)  & $-8.4\times 10^{-12}$ & $-0.01$ & 11/11
  & spreads \\
Concentrated jet & $+4.1\times 10^{-12}$ & $+0.92$ & 0/11
  & concentrates \\
\bottomrule
\end{tabular}
\end{table}

Three qualitatively different regimes emerge:

\medskip\noindent\textbf{Strong spreading (TG).}
The TG orbit shows $\Sigma \sim -10^{-3}$, orders of
magnitude larger than all other ICs, with perfect
anti-alignment $\cos\alpha=-1$ in high-$p$ sectors.
This is the regime where the $-i$ mechanism operates
at full strength, supported by the purely imaginary
phase structure.

\medskip\noindent\textbf{Noise-level (random broadband).}
The three random ICs have $|\Sigma|\sim 10^{-12}$,
nine orders of magnitude smaller than TG, and the sign
varies by seed (2/3 spread, 1/3 concentrate).  These
fields are nearly isotropic ($\cC\approx 0.07$, close
to $1/J=0.03$) and have conjugation-symmetric phases,
so $\E[\Sigma]=0$ by Theorem~\ref{thm:gaussian}.
The observed values are consistent with zero-mean
fluctuations.

\medskip\noindent\textbf{Concentration (jet).}
The concentrated jet shows $\Sigma > 0$ at every snapshot,
with concentration \emph{increasing} over time
($\cC: 0.475\to 0.478$).  Remarkably, the jet still
exhibits the $-i$ mechanism at the mode level:
$\corr(p_j,\langle\cos\alpha\rangle_j) = -0.70$
(the forcing depletes high-$p$ modes).  But $\Sigma>0$
despite this mode-level depletion.  The explanation:
the jet's random phases lack the cascade phase structure
needed for the cross-sector transfer to dominate.
The nonlinear term creates new angular structure
(at new scales and directions) faster than it depletes
the existing concentration.

\medskip\noindent\textbf{Degenerate (ABC).}
The ABC flow is a Beltrami field satisfying
$\nabla\times u=u$, making it a stationary Euler solution.
The nonlinear term $(u\cdot\nabla)u = (\nabla\times u)
\times u + \nabla(|u|^2/2)$ reduces to a pressure gradient
for Beltrami fields, so $B(u,u)\approx 0$ and $\Sigma$
is at machine zero.  This is not a counterexample but
a degenerate case where the nonlinearity barely acts.

\medskip
The jet result is the most informative: it shows that
the $-i$ mechanism (mode-level depletion) is necessary
but not sufficient for spreading.  Angular spreading
$\Sigma<0$ additionally requires the NS cascade phase
structure, which is present in developed TG turbulence
but absent in random-phase concentrated fields.


\subsection{Triad phase locking (N10)}
\label{sec:triad-phase}

The factored form (Proposition~\ref{prop:factored}) writes
the per-mode transfer as
$T(k) = 2|k|\,\im\bigl[\sum_{p+q=k}z(k;p,q)\bigr]$,
where the triad integrand is
\begin{equation}\label{eq:triad-integrand}
z(k;p,q) = (q\cdot\hat u(p))\,
(\hat u(k)^*\cdot\hat u(q)).
\end{equation}
We sample $2\times 10^5$ random triads from the set of
active modes (those with $|\hat u(k)|^2$ above $10^{-12}$
of the peak) and compute the magnitude-weighted phase
histogram of $\arg z$.

For the TG-generated NS solution at $t=10$,
$\nu=0.002$, $N=128$ (648 active modes,
$1.8\times 10^5$ nontrivial triads), the triad phases
concentrate sharply at $\arg z \approx +\pi/2$ (81\%
of total weight) and $\arg z\approx -\pi/2$ (18\%),
with all other bins essentially empty
(Table~\ref{tab:triad-phase}).  For comparison, the
random-phase field (same amplitudes) and the
concentrated jet both produce uniform histograms
indistinguishable from i.i.d.\ phases
($\approx 4\%$ per bin).

\begin{table}[h]
\centering
\caption{Triad phase statistics.
$\langle\sin\theta\rangle$ is the magnitude-weighted
mean of $\sin(\arg z)$; positive means net amplification
of total $\Hdot{1/2}$ energy by the nonlinear term.}
\label{tab:triad-phase}
\begin{tabular}{@{}lccc@{}}
\toprule
Field & triads & $\langle\sin\theta\rangle$
  & phase structure \\
\midrule
NS solution    & $1.8\times 10^5$ & $+0.63$
  & locked to $\pm\pi/2$ \\
Random phases  & $1.8\times 10^5$ & $-0.08$
  & uniform \\
Conc.\ jet     & $4.7\times 10^5$ & $-0.00$
  & uniform \\
\bottomrule
\end{tabular}
\end{table}

The phase locking has a transparent algebraic origin.
By Corollary~\ref{cor:persistence}, $\hat u(k)=iV(k)$
with $V(k)\in\R^3$ persists along the TG orbit.
Substituting into~\eqref{eq:triad-integrand}:
\begin{equation}\label{eq:triad-imaginary}
z(k;p,q) = (q\cdot iV(p))\,(-iV(k)\cdot iV(q))
= i\,(q\cdot V(p))(V(k)\cdot V(q)),
\end{equation}
which is purely imaginary, with
$\arg z = +\pi/2$ when $(q\cdot V(p))(V(k)\cdot V(q))>0$
and $\arg z = -\pi/2$ otherwise.  The observed $4\!:\!1$
bias toward $+\pi/2$ reflects the fact that the dominant
triads have both factors positive, so the aggregate
nonlinear term amplifies total $\Hdot{1/2}$ energy
($\langle\sin\theta\rangle>0$).

The key point for spreading is that this amplification
is \emph{anti-correlated with angular concentration}:
modes in low-$p_j$ sectors receive more nonlinear
amplification than those in high-$p_j$ sectors
(consistent with $\corr(p_j,dE_j|_{\rm NL})<0$
in~N1), producing $\Sigma<0$ even though the total
energy transfer is positive.

The uniform phase distribution for the random-phase
field and the concentrated jet confirms that the
$\pm\pi/2$ locking is a consequence of the
NS-preserved purely imaginary structure, not a
generic property of the trilinear form.  This
provides a positive characterization of the
``cascade phase structure'' invoked in N9 and
Conjecture~\ref{conj:enstrophy}: it is the persistence
of $\hat u(k)\approx iV(k)$, which locks triad
phases to the imaginary axis.


\subsection{Viscosity-independent constitutive relation (N11)}
\label{sec:constitutive}

We evolve the TG vortex from $t=0$ to $t=10$ at four
viscosities $\nu\in\{0.02, 0.01, 0.005, 0.002\}$
($N=128$, $J=32$), recording $\cC(u)$ and $\Sigma(u)$
at intervals $\Delta t=0.1$.

\begin{table}[h]
\centering
\caption{Mean spreading diagnostics over
$t\in[5,10]$.  The ratio $\Sigma/(\cC-1/J)$ is
independent of $\nu$ to four significant figures.}
\label{tab:constitutive}
\begin{tabular}{@{}cccc@{}}
\toprule
$\nu$ & $\langle\cC\rangle$
  & $\langle\Sigma/(\cC\!-\!1/J)\rangle$
  & std \\
\midrule
$0.02$  & $0.243$ & $-8.245\times 10^{-3}$ & $1.50\times 10^{-3}$ \\
$0.01$  & $0.243$ & $-8.256\times 10^{-3}$ & $1.50\times 10^{-3}$ \\
$0.005$ & $0.243$ & $-8.261\times 10^{-3}$ & $1.50\times 10^{-3}$ \\
$0.002$ & $0.243$ & $-8.264\times 10^{-3}$ & $1.51\times 10^{-3}$ \\
\bottomrule
\end{tabular}
\end{table}

The ratio $\Sigma/(\cC-1/J)$ agrees across all four
viscosities to four significant figures at every time
snapshot (Table~\ref{tab:constitutive}), confirming
that the spreading is purely nonlinear: viscous
dissipation contributes negligibly.  However, the ratio
evolves from $-2.2\times 10^{-3}$ at $t=2$ to
$-1.1\times 10^{-2}$ at $t=10$ (a factor of~5), while
$(\cC-1/J)$ itself decreases by only 5\%, indicating
that $\Sigma$ depends on an additional state variable
beyond~$\cC$.  Identifying the correct constitutive law
\begin{equation}\label{eq:constitutive}
\Sigma(u) = -g(u)\,(\cC(u)-1/J),
\end{equation}
where $g(u)>0$ is a viscosity-independent functional
of $u$, is a key open problem for the regularity
program.


\subsection{$\Hdot{s}$ spreading for $s\in[0,3/2]$ (N12)}
\label{sec:Hs-spreading}

\begin{definition}[$\Hdot{s}$ sector energies and spreading]
\label{def:Hs-spreading}
For $s\geq 0$, define the $\Hdot{s}$ sector energies
\begin{equation}
E_j^{(s)} = \sum_{\hat k\in\Omega_j}
|k|^{2s}\,|\hat u(k)|^2,
\qquad p_j^{(s)} = E_j^{(s)}/E^{(s)},
\qquad \cC^{(s)}(u) = \sum_j (p_j^{(s)})^2,
\end{equation}
and the $\Hdot{s}$ transfer and spreading functional
\begin{equation}
T^{(s)}(k;u) = 2|k|^{2s}\,\re[\hat u(k)^*\cdot
\hat B(k)],
\qquad
\Sigma^{(s)}(u) = \frac{d\cC^{(s)}}{dt}
\bigg|_{\rm NL}.
\end{equation}
The case $s=1/2$ recovers the original definitions.
For $s=1$, the sector energies $E_j^{(1)}$ are the
angular components of the enstrophy.
\end{definition}

We evolve the TG vortex at $N=128$, $\nu=0.005$
from $t=0$ to $t=10$, computing $\Sigma^{(s)}$ for
$s\in\{0,\,1/4,\,1/2,\,3/4,\,1,\,3/2\}$.

\begin{table}[h]
\centering
\caption{$\Hdot{s}$ spreading diagnostics, averaged
over $t\in[5,10]$, $J=32$.  The spreading strengthens
with increasing~$s$.}
\label{tab:Hs}
\begin{tabular}{@{}ccccl@{}}
\toprule
$s$ & $\langle\cC^{(s)}\rangle$
  & $\langle\Sigma^{(s)}\rangle$
  & $\langle\corr\rangle$
  & $\Sigma^{(s)}<0$\;always? \\
\midrule
$0$   & $0.245$ & $-1.11\times 10^{-3}$ & $-0.75$ & Yes \\
$1/4$ & $0.244$ & $-1.41\times 10^{-3}$ & $-0.68$ & Yes \\
$1/2$ & $0.243$ & $-1.79\times 10^{-3}$ & $-0.60$ & Yes \\
$3/4$ & $0.241$ & $-2.25\times 10^{-3}$ & $-0.51$ & Yes \\
$1$   & $0.238$ & $-2.82\times 10^{-3}$ & $-0.43$ & Yes \\
$3/2$ & $0.231$ & $-4.37\times 10^{-3}$ & $-0.30$ & Yes \\
\bottomrule
\end{tabular}
\end{table}

$\Sigma^{(s)}<0$ holds at every snapshot and every~$s$
(Table~\ref{tab:Hs}), and the spreading intensifies
with~$s$: $|\Sigma^{(3/2)}|/|\Sigma^{(1/2)}|\approx 2.4$.
The $s$-independence of the sign has a transparent
algebraic origin: the transfer
$T^{(s)}(k)=2|k|^{2s}\re[\hat u^*\cdot\hat B]$
shares the same phase structure as $T^{(1/2)}$
(the factor $|k|^{2s}$ is real and positive), so the
$\pm\pi/2$ locking of N10 drives
$\Sigma^{(s)}<0$ at every Sobolev level simultaneously.
The antisymmetry theorem (Theorem~\ref{thm:antisym}),
the vanishing theorem (Theorem~\ref{thm:gaussian}),
and the factored form (Proposition~\ref{prop:factored})
all extend verbatim to $\Sigma^{(s)}$.

The case $s=1$ is of particular significance for the
regularity problem: $\Sigma^{(1)}(u) < 0$ means the
nonlinearity spreads enstrophy angularly, opposing the
concentration that would be needed for vortex-sheet or
vortex-tube singularities.


\subsection{Vortex-stretching decomposition (N13)}
\label{sec:stretching}

The enstrophy equation for 3D Navier--Stokes reads
\begin{equation}\label{eq:enstrophy}
\frac{1}{2}\frac{d}{dt}\|\omega\|_{L^2}^2
= -\nu\|\nabla\omega\|_{L^2}^2
+ \int_{\T^3}\omega\cdot(\omega\cdot\nabla)u\,dx,
\end{equation}
where the vortex-stretching term
$\mathcal{W}=\int\omega\cdot(\omega\cdot\nabla)u\,dx$
is the sole mechanism for potential enstrophy blow-up.
We decompose $\mathcal{W}$ by angular sector.

Over $t\in[5,10]$ at $N=128$, $\nu=0.005$:
\begin{center}
\renewcommand{\arraystretch}{1.2}
\begin{tabular}{@{}lcc@{}}
\toprule
\textbf{Mechanism} & \textbf{corr with $p_j^{(1)}$}
  & \textbf{Effect} \\
\midrule
Full NL transfer $dE_j^{(1)}|_{\rm NL}$
  & $-0.43$ & net spreading \\
Vortex stretching $\mathcal{W}_j$
  & $+0.14$ & slight concentration \\
\bottomrule
\end{tabular}
\end{center}

The full nonlinear transfer produces net spreading
($\corr=-0.43$), while the stretching component alone
produces weak concentration ($\corr=+0.14$).  The
spreading is dominated by cross-sector transfer, which
overwhelms the concentrating tendency of stretching by
a factor of $\approx 3$ in correlation magnitude.
Along the TG orbit, $\cC^{(1)}$ is monotonically
decreasing: spreading wins.


\subsection{Viscosity independence at $s=1$ (N14)}
\label{sec:enst-visc}

We repeat the viscosity-independence test at the
enstrophy level.

\begin{table}[h]
\centering
\caption{Enstrophy-level constitutive relation.
$\Sigma^{(1)}/(\cC^{(1)}-1/J)$ is independent of
$\nu$ to four significant figures.}
\label{tab:enst-visc}
\begin{tabular}{@{}cccc@{}}
\toprule
$\nu$ & $\langle\cC^{(1)}\rangle$
  & $\langle\Sigma^{(1)}/(\cC^{(1)}\!-\!1/J)\rangle$
  & std \\
\midrule
$0.02$  & $0.238$ & $-1.367\times 10^{-2}$
  & $2.36\times 10^{-3}$ \\
$0.01$  & $0.238$ & $-1.369\times 10^{-2}$
  & $2.37\times 10^{-3}$ \\
$0.005$ & $0.238$ & $-1.370\times 10^{-2}$
  & $2.37\times 10^{-3}$ \\
$0.002$ & $0.238$ & $-1.370\times 10^{-2}$
  & $2.37\times 10^{-3}$ \\
\bottomrule
\end{tabular}
\end{table}

The constitutive relation
$\Sigma^{(1)}(u) = -g^{(1)}(u)\,(\cC^{(1)}(u)-1/J)$
holds with $g^{(1)}>0$ independent of $\nu$
(Table~\ref{tab:enst-visc}), confirming that
enstrophy-level spreading is entirely a property of
the nonlinear term.


\subsection{Multi-IC enstrophy spreading (N15)}
\label{sec:enst-multi-ic}

\begin{table}[h]
\centering
\caption{Enstrophy spreading ($s=1$) across initial
conditions at $N=128$, $\nu=0.01$, $J=32$.
Compare with Table~\ref{tab:multi-ic} ($s=1/2$).}
\label{tab:enst-multi-ic}
\begin{tabular}{@{}lcccl@{}}
\toprule
IC & $\Sigma^{(1)}(t\!=\!5)$ & $\corr$
  & $n_{\rm spread}$ & status \\
\midrule
Taylor--Green
  & $-1.93\times 10^{-3}$ & $-0.45$ & 10/10
  & spreads \\
ABC flow
  & $\approx 0$ & $+0.99$ & 2/10
  & degenerate \\
Random (seed 42)
  & $-4.2\times 10^{-11}$ & $-0.17$ & 10/10
  & spreads \\
Random (seed 99)
  & $-8.9\times 10^{-13}$ & $-0.00$ & 10/10
  & spreads \\
Random (seed 7)
  & $-1.6\times 10^{-11}$ & $-0.07$ & 10/10
  & spreads \\
Conc.\ jet
  & $+3.5\times 10^{-11}$ & $-0.53$ & 0/10
  & concentrates \\
\bottomrule
\end{tabular}
\end{table}

At $s=1/2$ (Table~\ref{tab:multi-ic}), the three
random broadband ICs gave mixed results: two seeds
spread and one concentrated.  At $s=1$
(Table~\ref{tab:enst-multi-ic}), \emph{all three
random seeds spread at every snapshot} (30/30 negative
values).  The probability of 30 consecutive negative
values from a distribution symmetric around zero
(Theorem~\ref{thm:gaussian}) is
$2^{-30}\approx 10^{-9}$.  Enstrophy-level spreading
is therefore more universal than $\Hdot{1/2}$ spreading.

The jet remains the only counterexample, with
$\Sigma^{(1)}>0$ at every snapshot, and the ABC flow
remains degenerate.


\subsection{Stretching landscape and the depletion
transition (N16)}
\label{sec:Phi-landscape}

To determine how vortex stretching depends on
angular concentration and phase structure jointly,
we construct synthetic fields by independently
modifying the concentration (sector amplitude
rescaling, parameter $\gamma$) and the phase
coherence (partial phase randomization, fraction $f$)
of a developed NS solution.

At full phase coherence ($f=0$), the mode-level
depletion metric
$\kappa_{\rm align} = -\sum_k w_k
\cos(\arg[\hat u(k)^*\cdot\hat B(k)])$
(enstrophy-weighted) undergoes a sharp sign
transition:
\begin{center}
\begin{tabular}{@{}ccc@{}}
\toprule
$\cC^{(1)}$ & $\kappa_{\rm align}$ & regime \\
\midrule
$0.04$ & $+1.00$ & amplification \\
$0.09$ & $+0.99$ & amplification \\
$0.15$ & $+0.83$ & transition \\
$0.21$ & $-0.56$ & depletion \\
$0.25$ & $-1.00$ & full depletion \\
\bottomrule
\end{tabular}
\end{center}

For nearly isotropic fields, the nonlinear term
\emph{amplifies} all modes: there is no dominant
sector to deplete.  As concentration increases past
$\cC^{(1)}\approx 0.15$, the cross-sector transfer
engages and the $-i$ mechanism begins depleting the
concentrated sectors.  This demonstrates that the
hypothesis ``$\cC^{(1)}>1/J$'' in the spreading
conjecture is structurally necessary.

A simple bound $\Phi\leq h(\cC^{(1)})$ with
$h(1/J)=0$ does not hold: the normalized stretching
$\Phi = |\mathcal{W}|/(\|\omega\|_{L^2}^2
\|\nabla u\|_{L^2})$ depends on the full phase
structure, not just angular concentration.
A multivariate regression of $\log\Phi$ against
$(\log(\cC^{(1)}-1/J),\,\kappa_{\rm imag})$ yields
$R^2=0.16$: the stretching landscape is not
well-described by simple scalar parameters.


\subsection{Spreading robustness under phase
destruction (N17)}
\label{sec:robustness}

\begin{center}
\begin{tabular}{@{}cccc@{}}
\toprule
$f$ & $\Sigma^{(1)}<0$ & $\langle\Sigma^{(1)}\rangle$
  & $\langle\kappa_{\rm align}\rangle$ \\
\midrule
$0.0$ & 11/11 & $-1.74\times 10^{-3}$ & $-0.43$ \\
$0.1$ & 10/11 & $-1.27\times 10^{-3}$ & $-0.36$ \\
$0.3$ & 10/11 & $-8.30\times 10^{-4}$ & $-0.30$ \\
$0.5$ & 10/11 & $-3.04\times 10^{-4}$ & $-0.04$ \\
$0.7$ &  5/11 & $-1.93\times 10^{-4}$ & $-0.04$ \\
$1.0$ &  7/11 & $+9.25\times 10^{-6}$ & $+0.10$ \\
\bottomrule
\end{tabular}
\end{center}

Enstrophy spreading persists at nearly all
concentration levels ($\geq 10/11$) up to 50\%
phase randomization.  The breakdown around $f=0.7$
coincides with $\kappa_{\rm align}$ crossing zero.
Combined with N15 (where random broadband ICs develop
$\Sigma^{(1)}<0$ under NS evolution), this suggests
the NS dynamics drive sufficient phase coherence
for spreading even from incoherent initial data.


\subsection{Surgical concentration and self-regulation (N18)}
\label{sec:surgical}

The preceding tests evolve NS solutions from prescribed
initial conditions and measure spreading along the
resulting trajectories.  A stronger test of the
regularity program is to \emph{artificially concentrate}
a developed NS solution---preserving its phase
structure---and ask whether the dynamics restore
isotropy.

\medskip\noindent\textbf{Construction.}
We evolve the TG vortex at $N=256$, $\nu=0.01$ to
$t=5$, producing a developed turbulent field with
$\cC^{(1)}\approx 0.042$.  We then surgically
increase enstrophy concentration by amplifying the
dominant sector by a factor $\gamma>1$ and suppressing
all others by $\gamma^{-1/2}$, with $\gamma$ chosen
by bisection to reach target concentrations
$\cC^{(1)}\in\{0.20, 0.30, 0.40, 0.50\}$.  The
modified field is Leray-projected to restore
divergence-freeness.  Crucially, this procedure
preserves the Fourier phase structure generated by
NS evolution: only the sector amplitudes change.
Each surgically concentrated field is then evolved
under NS for $t\in[0,5]$ at the same $\nu=0.01$.

\medskip\noindent\textbf{Results.}
At all 44 diagnostic snapshots across all four
concentration levels, $\Sigma^{(1)}<0$ without
exception.  Enstrophy concentration decreases
monotonically from its surgically imposed value:

\begin{table}[h]
\centering
\caption{Surgical concentration: NS relaxation at
$N=256$, $\nu=0.01$, $J=32$.  $\Sigma^{(1)}<0$ at
every snapshot (44/44).}
\label{tab:surgical}
\begin{tabular}{@{}cccccc@{}}
\toprule
$\cC^{(1)}(0)$ & $\cC^{(1)}(5)$
  & peak $|\Sigma^{(1)}|$
  & $|W|$ range
  & $\mathcal{R}$ range
  & $g^{(1)}$ range \\
\midrule
$0.20$ & $0.13$ & $7.95\times 10^{-2}$
  & $11$--$34$ & $1.8$--$2.2$ & $0.40$--$0.56$ \\
$0.30$ & $0.17$ & $1.16\times 10^{-1}$
  & $11$--$30$ & $1.1$--$1.4$ & $0.33$--$0.52$ \\
$0.40$ & $0.21$ & $1.47\times 10^{-1}$
  & $11$--$29$ & $0.8$--$1.0$ & $0.27$--$0.49$ \\
$0.50$ & $0.23$ & $1.74\times 10^{-1}$
  & $12$--$27$ & $0.6$--$0.8$ & $0.22$--$0.48$ \\
\midrule
$0.042$ (base) & $0.040$ & $2.87\times 10^{-3}$
  & $17$--$78$ & $68$--$295$ & $0.08$--$0.31$ \\
\bottomrule
\end{tabular}
\end{table}

Three features of the data are significant:

\medskip\noindent\textbf{(i) Monotone relaxation.}
Enstrophy concentration $\cC^{(1)}(t)$ decreases at
every snapshot in every surgical run.  The NS
nonlinearity does not permit sustained angular
concentration: when concentration is artificially
elevated, the dynamics pull it back toward isotropy
with a rate proportional to the excess
$\cC^{(1)}-1/J$.

\medskip\noindent\textbf{(ii) Spreading scales with
concentration; stretching does not.}
The peak spreading rate $|\Sigma^{(1)}|$ increases
from $0.080$ to $0.174$ as the initial concentration
rises from $0.20$ to $0.50$ (a factor of $2.2\times$).
By contrast, vortex stretching $|W|$ remains in the
range $11$--$34$ across all surgical runs---essentially
the same range as the unmodified base trajectory.
This decoupling is the quantitative content of
self-regulation: the spreading mechanism responds
to concentration, while the stretching mechanism
does not.

\medskip\noindent\textbf{(iii) The ratio
$\mathcal{R}$ decreases with concentration.}
Define the stretching--spreading ratio
\begin{equation}\label{eq:R-ratio}
\mathcal{R}(t) = \frac{|\mathcal{W}(t)|}
{|\Sigma^{(1)}(t)|\cdot\|\omega(t)\|_{L^2}^2},
\end{equation}
which measures how efficiently the nonlinear
spreading functional $\Sigma^{(1)}$ counteracts
vortex stretching $\mathcal{W}$, normalized by
enstrophy.  In the surgical runs,
$\mathcal{R}\in[0.6, 2.2]$ and \emph{decreases}
as concentration increases
(Table~\ref{tab:surgical}).  This is the opposite
of the blow-up scenario: the higher the
concentration, the more effectively spreading
dominates stretching.

The large values $\mathcal{R}\approx 70$--$295$
in the base run are harmless: they occur at
$\cC^{(1)}\approx 0.04 \approx 1/J$, where the
field is already isotropic and $|\Sigma^{(1)}|$ is
small because there is no angular concentration to
spread.  In this regime, neither stretching nor
spreading drives the dynamics: viscous dissipation
dominates.


\subsection{The constitutive law (N19)}
\label{sec:constitutive-surgical}

Sections~\ref{sec:constitutive}
and~\ref{sec:enst-visc} established the constitutive
relation $\Sigma^{(1)} = -g^{(1)}(u)\,
(\cC^{(1)}-1/J)$ along single TG trajectories, where
$g^{(1)}$ is viscosity-independent but varies in
time by a factor of~5.  The earlier regression of
$g^{(1)}$ against scalar parameters
(Section~\ref{sec:Phi-landscape}) gave $R^2=0.16$.
The surgical data resolve this.

\medskip\noindent\textbf{A nearly universal
constant.}
By pooling all snapshots from the four surgical runs
at $t\geq 1.5$ (after the initial transient), the
constitutive coefficient
\begin{equation}\label{eq:g1-value}
g^{(1)} = \frac{-\Sigma^{(1)}}{\cC^{(1)}-1/J}
\end{equation}
stabilizes at
\begin{equation}\label{eq:g1-numerical}
g^{(1)} \approx 0.47 \pm 0.04
\end{equation}
across the full range
$\cC^{(1)}\in[0.13, 0.50]$
(Table~\ref{tab:surgical}).  The coefficient of
variation is under $10\%$.

\begin{remark}[Why the single-trajectory regression
failed]
Along a single TG trajectory, $\cC^{(1)}$ and the
cascade age (measured by $k_*$ or $t$) are perfectly
correlated: both evolve together.  The regression
cannot distinguish their effects.  The surgical
construction decouples them: all runs share the same
cascade age (turbulence developed to $t=5$) but have
different concentrations.  The resulting near-constant
$g^{(1)}\approx 0.47$ shows that the constitutive
coefficient depends on the cascade phase structure
(which is shared), not on $\cC^{(1)}$ itself.
\end{remark}

\begin{remark}[Analytical target]
\label{rem:analytical-target}
The numerical constitutive law
\begin{equation}\label{eq:constitutive-target}
\Sigma^{(1)}(u) \approx -0.47\,
(\cC^{(1)}(u)-1/J)
\end{equation}
reduces the regularity program to a single
analytical question: prove that
$g^{(1)}(u) \geq c > 0$ for NS solutions with
developed cascade phase structure.  The factored
form (Proposition~\ref{prop:factored}) writes both
$\Sigma^{(1)}$ and $\mathcal{W}$ as sums over the
same triads $\{(k,p,q): p+q=k\}$.  The triad phase
locking of N10 constrains each summand to the
imaginary axis.  Bounding $g^{(1)}$ from below is
equivalent to bounding the ratio of two sums over
the same triads with the same phase structure---a
problem in the combinatorics of lattice convolutions
weighted by $|k|^2$, not a PDE estimate.
\end{remark}

\begin{remark}[Connection to Step~(ii)]
\label{rem:step-ii}
The boundedness of $\mathcal{R}$ in the surgical
runs (Table~\ref{tab:surgical}) is a direct
numerical verification of Step~(ii) in
Remark~\ref{rem:regularity-final}: the
stretching--spreading ratio remains bounded along
NS trajectories, even under adversarial
concentration.  Combined with Step~(i) (the
enstrophy spreading conjecture, supported by
N12--N17 and the 44/44 surgical snapshots), this
closes the numerical side of the regularity program.
The remaining challenge is analytical: proving
$g^{(1)}\geq c > 0$ from the trilinear algebra.
\end{remark}


%======================================================================
\section{A Heuristic Mechanism: the $-i$ Rotation}
\label{sec:mechanism}
%======================================================================

The exact results of Sections~\ref{sec:general}--\ref{sec:TG}
and the numerical findings of Section~\ref{sec:results}
suggest a general mechanism.  \textbf{The arguments in this
section and the next are heuristic: they rely on
concentration and slow-variation approximations that are
not rigorously justified for general NS solutions.}
We present them as a proposed explanation, not as proved
theorems.

\subsection{Polarization}

\begin{definition}[Polarization]
\label{def:polarization}
For a mode $\hat u(k)\neq 0$, the polarization is
$\hat e(k) = \hat u(k)/|\hat u(k)|\in\C^3$.  When $u$ is
angularly concentrated in sector $\Omega_{j^*}$ with
slowly varying polarization, we write
$\hat u(q)\approx A(q)\hat e$ for
$\hat q\in\Omega_{j^*}$, where $A(q)=|\hat u(q)|$ and
$\hat e$ is the (approximately) common polarization
direction.  For a cross-sector mode
$p\notin\Omega_{j^*}$, we write
$\hat e_p = \hat u(p)/|\hat u(p)|$.
\end{definition}

\subsection{Cross-sector forcing under concentration}

Under the concentration hypothesis, the cross-sector
forcing at a mode $k\in\Omega_{j^*}$ takes the
approximate form
\begin{equation}\label{eq:Bcross-approx}
\hat B^{\rm cross}(k) \approx -iF_{\rm cross}(k)\hat e,
\end{equation}
where
\begin{equation}
F_{\rm cross}(k) = \sum_{\substack{p+q=k\\
\hat p\notin\Omega_{j^*},\;\hat q\in\Omega_{j^*}}}
(q\cdot\hat u(p))\,A(q).
\end{equation}
The key step: since $\hat u(k)\approx A_k\hat e$ and
$P(k)\hat e\approx\hat e$ (both from the narrow-sector
hypothesis), the transfer becomes
$T^{\rm cross}(k) \approx 2|k|A_k\,\im[F_{\rm cross}(k)]$.

Spreading of the dominant sector therefore requires
$\im[F_{\rm cross}(k)] < 0$ for most modes
$k\in\Omega_{j^*}$.

\subsection{Why the TG case is exact}

For TG, the concentration and narrow-sector hypotheses
are not approximations: the 8 modes have identical
amplitudes and the purely imaginary structure makes
$P(k)\hat e = \hat e$ exact (since $V(k)\perp k$
for every $k\in\cK$).  The approximate mechanism becomes
exact, and $\im[F_{\rm cross}]<0$ is established by
Theorem~\ref{thm:TG}.

\subsection{The $-i$ rotation in general}

The factor $-i$ in the advection operator
$\widehat{(u\cdot\nabla)u}(k) = \sum_{p+q=k}
iq_l\hat u_l(p)\hat u(q)$ creates a $\pi/2$ phase
rotation between the ``natural'' phase of a triadic
interaction and the resulting forcing.  For concentrated
fields, this means cross-sector forcing (where the
advecting velocity comes from a different angular
direction) arrives with a systematic phase offset
relative to the existing velocity in the concentrated
sector.  The offset biases $\cos\alpha$ toward negative
values (depletion rather than amplification).

For TG, the offset is exactly $\pi$ (perfect
anti-alignment), because the purely imaginary structure
locks the accumulated phase to $i^2=-1$.  For general
fields, the offset is partial, producing
$\cos\alpha\approx-0.6$ on average in high-$p$ sectors
(N8).

The triad phase locking observed in N10 provides
direct microscopic evidence for this mechanism.
Equation~\eqref{eq:triad-imaginary} shows that the
purely imaginary structure forces every individual
triad integrand $z(k;p,q)$ onto the imaginary axis,
so the nonlinear transfer $T(k) = 2|k|\,\im[\sum z]$
is entirely controlled by the signs of the real
products $(q\cdot V(p))(V(k)\cdot V(q))$.  For
random-phase fields, these products are complex
with uniformly distributed arguments, and the
$\pm\pi/2$ locking disappears.  The cascade phase
structure required by Conjecture~\ref{conj:enstrophy}
is therefore precisely the persistence of
$\hat u\approx iV$, which the NS dynamics maintain
but arbitrary initial data need not possess.

The extension to $\Sigma^{(s)}$ for general~$s$ (N12)
follows because the factor $|k|^{2s}$ in $T^{(s)}(k)$
is a positive real weight that preserves the triad
phase structure.  The $\pm\pi/2$ locking drives
spreading at every Sobolev level simultaneously: the
nonlinear term redistributes the \emph{entire regularity
structure} of the field --- including the enstrophy
($s=1$) and higher derivatives ($s=3/2$) --- away from
angular concentration.

The depletion transition of N16 reveals a threshold
effect: the $-i$ mechanism engages only when angular
concentration exceeds $\cC^{(1)}\approx 0.15$.  Below
this threshold, the same NS phases amplify rather than
deplete, because there is no dominant sector for
cross-sector transfer to target.  The spreading
robustness of N17 shows that 50\% of the phase
structure can be destroyed before the mechanism fails,
establishing that the $-i$ rotation is not fragile.

%======================================================================
\section{A Conditional Sign Criterion}
\label{sec:conditional}
%======================================================================

\textbf{This section presents a conditional result:
depletion under an explicit sign hypothesis.  It does
not constitute a rigorous reduction to Kolmogorov's
forward cascade.}

\subsection{Quasi-static cross-sector modes}

If cross-sector modes equilibrate viscously
($\nu|k'|^2$ large compared to the forcing variation
rate), the NS equation gives
$\hat u(k')\approx\hat B(k')/(\nu|k'|^2)$.  Since
$\hat B(k')\approx -iG(k')$ for modes created by
within-sector triads (with
$G(k')=P(k')\sum_{p,q\in\Omega_{j^*}}
(q\cdot\hat u(p))\hat u(q)$), the cross-sector mode
carries a phase approximately $\pi/2$ behind its source.

This approximation is rigorous in the dissipation range
where $\nu|k'|^2\gg\tau_{\rm NL}^{-1}$.  In the
inertial range, where the nonlinear turnover time
competes with viscous damping, the quasi-static
assumption softens, but the $\pi/2$ bias from the $-i$
factor persists as a leading-order effect.

\subsection{The sign criterion}

When quasi-static cross-sector modes feed back to
$k\in\Omega_{j^*}$, the feedback transfer involves
\begin{equation}
\im[\gamma(k,p')] = \im\Bigl[q\cdot P(p')
\sum_{\substack{r+s=p'\\r,s\in\Omega_{j^*}}}
(s\cdot\hat u(r))\hat u(s)\Bigr],
\end{equation}
where $q=k-p'$.  If this quantity is negative for
the dominant intermediate wavenumbers $p'$, the feedback
depletes the concentrated sector.

\begin{remark}[Not a reduction to K41]
We emphasize that $\im[\gamma]<0$ is a sign condition
on the same trilinear functional whose negativity
one wants to establish.  It is \emph{not} equivalent
to the standard forward cascade (constant energy flux
through spectral shells).  The connection to Kolmogorov's
theory is suggestive --- both involve the sign of
triadic phase interactions --- but the condition here is
sector-projected and geometry-dependent, not a direct
consequence of K41.  We state it as a conditional
criterion, not as a proved reduction.
\end{remark}

%======================================================================
\section{Conjecture and Discussion}
\label{sec:conjecture}
%======================================================================

\begin{conjecture}[Enstrophy spreading for developed
turbulence]
\label{conj:enstrophy}
For smooth NS solutions on $\T^3$ that have developed
cascade phase structure (in the sense of having evolved
from generic initial data for a time sufficient to
establish a forward energy cascade): if
$\cC^{(1)}(u(t))>1/J$, then $\Sigma^{(1)}(u(t))<0$.
\end{conjecture}

The evidence for this conjecture consists of:
the exact algebraic extension of N12 (the
$|k|^{2s}$ weight preserves the triad phase
structure for all~$s$);
the TG result (N12: $\Sigma^{(1)}<0$ at every
snapshot);
the viscosity independence (N14);
the multi-IC evidence (N15: all random broadband
ICs spread at $s=1$, 30/30 negative snapshots);
the jet counterexample (N15: cascade phase structure
remains necessary);
the depletion transition (N16: the mechanism engages
above $\cC^{(1)}\approx 0.15$);
and the phase robustness (N17: spreading survives
50\% phase destruction);
the surgical concentration test (N18: $\Sigma^{(1)}<0$
at all 44 snapshots under adversarial concentration up
to $\cC^{(1)}=0.50$, with monotone relaxation);
and the constitutive law (N19: $g^{(1)}\approx
0.47\pm 0.04$ across the full concentration range).

Conjecture~\ref{conj:enstrophy} implies the $s=1/2$
version, since the triad phase structure that drives
$\Sigma^{(1)}<0$ simultaneously drives
$\Sigma^{(s)}<0$ for all $s\geq 0$ (N12).

\begin{conjecture}[$\Hdot{1/2}$ spreading]
\label{conj:main}
Under the same hypotheses: if $\cC(u(t))>1/J$,
then $\Sigma(u(t))<0$.
\end{conjecture}

\begin{remark}[Connection to regularity]
\label{rem:regularity-final}
The regularity question reduces to two coupled
statements about the NS nonlinearity.

\smallskip\noindent\textbf{Step (i): Enstrophy
spreading.}
Conjecture~\ref{conj:enstrophy}:
$\Sigma^{(1)}(u)<0$ whenever $\cC^{(1)}>1/J$.
Evidence: N12 (all Sobolev levels), N14 (viscosity
independence), N15 (generality), N17 (robustness).

\smallskip\noindent\textbf{Step (ii):
Stretching--spreading coupling.}
The surgical concentration tests of N18 provide
direct evidence: under adversarial enstrophy
concentration ($\cC^{(1)}$ up to $0.50$), the ratio
$\mathcal{R} = |\mathcal{W}|/(|\Sigma^{(1)}|
\cdot\|\omega\|^2)$ remains in $[0.6, 2.2]$ and
\emph{decreases} with concentration
(Table~\ref{tab:surgical}).  The constitutive law
$g^{(1)}\approx 0.47$ (N19) identifies the remaining
analytical challenge: prove $g^{(1)}\geq c > 0$ for NS
solutions with developed cascade phase structure.
The factored form (Proposition~\ref{prop:factored})
reduces this to a combinatorial bound on lattice
triads weighted by $|k|^2$.

Formalizing this bound --- proving that the
stretching--spreading ratio $\mathcal{W}/|\Sigma^{(1)}|$
remains bounded along NS trajectories --- is the
remaining analytical challenge.
\end{remark}

The elevation of the conjecture from $s=1/2$ to $s=1$
has a structural consequence.  At $s=1/2$, the
connection to regularity required an intermediate
``angular Bernstein inequality'' relating concentration
of $\Hdot{1/2}$ energy to pointwise norms.  At $s=1$,
this intermediate step is bypassed: enstrophy
concentration is already the regularity-relevant
quantity through the vortex-stretching term
in~\eqref{eq:enstrophy} and the Beale--Kato--Majda
criterion~\cite{BKM}.  The regularity problem
reduces to questions about the trilinear form and its
angular geometry --- precisely the domain where the
algebraic identities of
Sections~\ref{sec:general}--\ref{sec:TG} provide
leverage.

%======================================================================
\section{Summary}
%======================================================================

\begin{center}
\renewcommand{\arraystretch}{1.2}
\begin{tabular}{@{}lll@{}}
\toprule
\textbf{Result} & \textbf{Status} & \textbf{Section} \\
\midrule
$T(k;\bar u) = -T(k;u)$
  & Proved & \ref{sec:general} \\
$\E_{\bar u=u}[T(k)]=0$, $\E[\Sigma]=0$
  & Proved & \ref{sec:general} \\
$T(k)$ is pressure-invariant
  & Proved & \ref{sec:general} \\
$\Sigma(\lambda u)=\lambda\Sigma(u)$
  & Proved & \ref{sec:general} \\
$V(k)\cdot\widetilde R(k)=-1/512$ (TG)
  & Proved & \ref{sec:TG} \\
$\Sigma<0$ along TG orbit
  & Numerical & \ref{sec:results} \\
$\nu$-independence of spreading
  & Numerical & \ref{sec:results} \\
Cross-sector dominance ($830\times$)
  & Numerical & \ref{sec:results} \\
Jet concentrates despite mode-level depletion
  & Numerical & \ref{sec:multi-ic} \\
Triad phases locked to $\pm\pi/2$ (NS only)
  & Numerical & \ref{sec:triad-phase} \\
$\Sigma/(\cC-1/J)$ independent of $\nu$
  & Numerical & \ref{sec:constitutive} \\
$\Sigma^{(s)}<0$ for $s\in[0,3/2]$
  & Numerical & \ref{sec:Hs-spreading} \\
Stretching vs.\ spreading ($+0.14$ vs.\ $-0.43$)
  & Numerical & \ref{sec:stretching} \\
$\Sigma^{(1)}/(\cC^{(1)}-1/J)$ indep.\ of $\nu$
  & Numerical & \ref{sec:enst-visc} \\
$\Sigma^{(1)}<0$ for TG + all 3 random ICs
  & Numerical & \ref{sec:enst-multi-ic} \\
$\kappa_{\rm align}$ sign transition at $\cC^{(1)}\approx 0.15$
  & Numerical & \ref{sec:Phi-landscape} \\
$\Sigma^{(1)}$ robust to 50\% phase randomization
  & Numerical & \ref{sec:robustness} \\
Surgical: $\Sigma^{(1)}<0$ at all 44 snapshots
  & Numerical & \ref{sec:surgical} \\
$\mathcal{R}\in[0.6,2.2]$ under adversarial concentration
  & Numerical & \ref{sec:surgical} \\
$g^{(1)}\approx 0.47\pm 0.04$ (constitutive law)
  & Numerical & \ref{sec:constitutive-surgical} \\
$-i$ rotation mechanism
  & Heuristic & \ref{sec:mechanism} \\
$\im[\gamma]<0$ sign criterion
  & Conditional & \ref{sec:conditional} \\
Conjecture~\ref{conj:enstrophy} ($s=1$ spreading)
  & Conjecture & \ref{sec:conjecture} \\
Stretching--spreading coupling ($\mathcal{R}$ bounded)
  & Numerical & \ref{sec:surgical} \\
$g^{(1)}\geq c > 0$ (analytical)
  & Open & \ref{sec:constitutive-surgical} \\
\bottomrule
\end{tabular}
\end{center}

\subsection*{Acknowledgments}

AI tools were used as interactive assistants for drafting
and mathematical exploration.
Simulations were performed on an NVIDIA RTX~5090.

%======================================================================
\begin{thebibliography}{99}

\bibitem{Fefferman}
C.\,L.~Fefferman,
Existence and smoothness of the Navier--Stokes equation,
\emph{Clay Mathematics Institute Millennium Prize Problems}
(2000), 1--5.

\bibitem{Leray}
J.~Leray,
Sur le mouvement d'un liquide visqueux emplissant l'espace,
\emph{Acta Math.}~\textbf{63} (1934), 193--248.

\bibitem{ESS}
L.~Escauriaza, G.\,A.~Seregin, V.~\v{S}ver\'ak,
$L_{3,\infty}$-solutions of Navier--Stokes equations and backward
uniqueness, \emph{Russian Math.\ Surveys}~\textbf{58} (2003),
211--250.

\bibitem{BKM}
J.\,T.~Beale, T.~Kato, A.~Majda,
Remarks on the breakdown of smooth solutions for the 3D Euler
equations, \emph{Comm.\ Math.\ Phys.}~\textbf{94} (1984), 61--66.

\bibitem{Tao}
T.~Tao,
Finite time blowup for an averaged three-dimensional
Navier--Stokes equation,
\emph{J.\ Amer.\ Math.\ Soc.}~\textbf{29} (2016), 601--674.

\bibitem{Kolmogorov}
A.\,N.~Kolmogorov,
The local structure of turbulence in incompressible viscous fluid
for very large Reynolds numbers,
\emph{Dokl.\ Akad.\ Nauk SSSR}~\textbf{30} (1941), 299--303.

\bibitem{Rotta}
J.\,C.~Rotta,
Statistische Theorie nichthomogener Turbulenz,
\emph{Z.\ Physik}~\textbf{129} (1951), 547--572.

\bibitem{CF93}
P.~Constantin and C.~Fefferman,
Direction of vorticity and the problem of global
regularity for the Navier--Stokes equations,
\emph{Indiana Univ.\ Math.\ J.}~\textbf{42} (1993),
775--789.

\end{thebibliography}

\end{document}
